\section{Coordinates and partial traces of party products}
\label{sec:peps_party_matrix_coordinates}

\begin{definition}[Party tensor basis]
    \label{def:peps_party_tensor_basis}
    \lean{TNLean.PEPS.PairEffect.partyListBasis,
      TNLean.PEPS.PairEffect.partyLabelledBasis}
    \leanok
    For each party $p$, choose an orthonormal basis $b_p$ of its input
    memory and an orthonormal basis $c_p$ of its output memory. An
    exhaustive party order without repetitions determines their ordered
    tensor bases. Relabel the coordinates by original parties, so that
    a basis vector is indexed by a function $p\mapsto i_p$.
    The empty tensor basis is the scalar unit.
\end{definition}

\begin{theorem}[Entries of the local tensor product]
    \label{thm:peps_party_tensor_entries}
    \lean{TNLean.PEPS.PairEffect.partyListBasis_repr_tensorPartyMaps,
      TNLean.PEPS.PairEffect.partyLabelledBasis_repr_tensorPartyMaps}
    \leanok
    \uses{def:peps_party_tensor_basis}
    Let $B_p$ be a linear map between the chosen input and output
    memories. In the party tensor bases, the matrix entries satisfy
    \begin{align}
        \left(\bigotimes_p B_p\right)_{x,y}
        &=\prod_p\langle c_p(x_p),B_pb_p(y_p)\rangle.
        \label{eq:peps_party_tensor_entries}
    \end{align}
    No bound on the local dimensions or map norms is needed.
\end{theorem}

\begin{proof}\leanok
    \uses{def:peps_party_tensor_basis}
    Induct on the party order. The coordinate of a simple tensor is
    the product of its two coordinates. Apply the first local map,
    use the induction hypothesis on the remaining maps, and separate
    the first factor of the finite product. Relabelling the exhaustive
    order replaces the product over positions by the product over parties.
\end{proof}

\begin{theorem}[Private trace of a product matrix]
    \label{thm:peps_party_product_trace}
    \lean{TNLean.PEPS.PairEffect.partialTraceRight_partyProductMatrix}
    \leanok
    Let $P$ be finite. For each $p\in P$, let $X_p$ be a physical
    coordinate set, let $E_p$ be a finite private coordinate set, and
    let $M_p$ be a complex matrix indexed by $X_p\times E_p$.
    Define $M_{(x,i),(y,j)}=\prod_p(M_p)_{(x_p,i_p),(y_p,j_p)}$.
    Then
    \begin{align}
        \operatorname{Tr}_{\prod_p E_p}M
        &=\bigotimes_p\operatorname{Tr}_{E_p}M_p.
        \label{eq:peps_party_product_trace}
    \end{align}
    The matrices may be arbitrary; positivity is not assumed.
\end{theorem}

\begin{proof}\leanok
    At coordinates $x,y$, the left side is
    $\sum_i\prod_p(M_p)_{(x_p,i_p),(y_p,i_p)}$.
    Finite distributivity makes this
    $\prod_p\sum_{i_p}(M_p)_{(x_p,i_p),(y_p,i_p)}$, which is the
    corresponding entry on the right.
\end{proof}

\begin{theorem}[Local conjugation followed by the private trace]
    \label{thm:peps_party_product_conjugation}
    \lean{TNLean.PEPS.PairEffect.partyProductMatrix_conjTranspose,
      TNLean.PEPS.PairEffect.partyProductMatrix_conjugate,
      TNLean.PEPS.PairEffect.partialTraceRight_partyProduct_conjugate}
    \leanok
    \uses{thm:peps_party_product_trace}
    Let $X_p$ be an arbitrary matrix on the source coordinates at party
    $p$, and let $A_p,B_p$ map these coordinates to the physical and
    private output coordinates at that party. Then
    \begin{align}
        \operatorname{Tr}_{\rm private}
          \left[(\bigotimes_p A_p)(\bigotimes_p X_p)
          (\bigotimes_p B_p)^*\right]
        &=\bigotimes_p\operatorname{Tr}_{\rm private}
          (A_pX_pB_p^*).
    \end{align}
\end{theorem}

\begin{proof}\leanok
    \uses{thm:peps_party_product_trace}
    Taking the adjoint takes the adjoint of each factor. Finite
    distributivity computes both products separately at each party.
    Apply the product partial-trace identity.
\end{proof}

These coordinate and trace identities express the local contractions in
\cite[proof of Theorem~5.2]{OpenAI2026PolynomialPEPS},
\texttt{04-compression.tex}, lines 565--579.
